\documentclass[
    fontsize=11pt,
    letterpaper,oneside,paper=portrait
]{scrartcl}

\usepackage[charter]{mathdesign}
\usepackage[sups,osf]{XCharter}
\usepackage[semibold]{sourcesanspro}
\usepackage[scaled=1.04,varqu,varl]{inconsolata}

\usepackage{microtype}
\usepackage[x11names,svgnames,dvipsnames,table]{xcolor}
\usepackage{hyperref}
\usepackage{booktabs}
\usepackage{array}
\usepackage{nicefrac}
\usepackage{longtable}

\usepackage[margin=0.75in]{geometry}

%%%% METADATA  %%%% %%%% %%%% %%%%
\newcommand{\myname}{lipidless.com}
\newcommand{\mydate}{29~April 2024}
\newcommand{\mytitle}{German potato salad}
\newcommand{\myurl}{https://www.lipidless.com/recipes/german-potato-salad}
%%%% %%%% %%%% %%%% %%%% %%%% %%%%

% links
\hypersetup{
  colorlinks=false,
  pdfborderstyle={/S/S/W 0.5},
  pdfborder={0 0 0.5},
  allbordercolors={0.8 0.8 1}
}

% no section numbering
\setcounter{secnumdepth}{0}

% reduce paragraph vspace
\KOMAoptions{parskip=half}

% do not indent paragraphs
\setlength{\parindent}{0pt}

% section heading styling
\addtokomafont{section}{\sffamily\Large}
\addtokomafont{subsection}{\sffamily\large}

% bold sans
\newcommand{\sfbf}[1]{\textbf{\sffamily #1}}

% bold sans with some spacing, like a compact paragraph
\newcommand{\paragrcompact}[1]{\textbf{\sffamily #1}\quad}

% shared recipe layout lengths
\newlength{\recipeamountwidth}
\newlength{\recipeunitwidth}
\newlength{\recipedescnaturalwidth}
\newlength{\recipeingredientdescwidth}
\newlength{\recipeingredientblockwidth}
\newlength{\recipestepwidth}
\newlength{\recipecolumngap}
\newlength{\recipeamountunitgap}
\newlength{\recipeunitdescgap}
\newlength{\recipewidthbuffer}
\newcommand{\measuremaxwidth}[2]{

  \settowidth{\dimen0}{#2}%
  \ifdim\dimen0>#1\relax
    #1=\dimen0\relax
  \fi
}

\pagestyle{empty}
\setlength{\LTleft}{0pt}
\setlength{\LTright}{0pt}

\begin{document}

{\sffamily \textbf{\LARGE \mytitle} \hfill \mydate \hfill \href{\myurl}{\myname}}

\vspace{\baselineskip}

Cold potato salad using mustard and vinegar (no mayonnaise). Potatoes freeze poorly, so this is best kept refrigerated.

\subsection*{Steps}

\renewcommand{\arraystretch}{1.15}
\setlength{\recipecolumngap}{1.0em}
\setlength{\recipeamountunitgap}{0.3em}
\setlength{\recipeunitdescgap}{0.8em}
\setlength{\recipewidthbuffer}{0.35em}
\setlength{\recipeamountwidth}{2em}
\setlength{\recipeunitwidth}{0pt}
\setlength{\recipedescnaturalwidth}{0pt}
\measuremaxwidth{\recipeamountwidth}{5}
\measuremaxwidth{\recipeunitwidth}{lbs}
\measuremaxwidth{\recipedescnaturalwidth}{medium russet potatoes}
\measuremaxwidth{\recipeamountwidth}{1}
\measuremaxwidth{\recipeunitwidth}{c}
\measuremaxwidth{\recipedescnaturalwidth}{water}
\measuremaxwidth{\recipeamountwidth}{2}
\measuremaxwidth{\recipeunitwidth}{\mbox{}}
\measuremaxwidth{\recipedescnaturalwidth}{chopped yellow onions}
\measuremaxwidth{\recipeamountwidth}{1.5}
\measuremaxwidth{\recipeunitwidth}{c}
\measuremaxwidth{\recipedescnaturalwidth}{water}
\measuremaxwidth{\recipeamountwidth}{1/2}
\measuremaxwidth{\recipeunitwidth}{c}
\measuremaxwidth{\recipedescnaturalwidth}{vinegar}
\measuremaxwidth{\recipeamountwidth}{2}
\measuremaxwidth{\recipeunitwidth}{t}
\measuremaxwidth{\recipedescnaturalwidth}{sea salt}
\measuremaxwidth{\recipeamountwidth}{1}
\measuremaxwidth{\recipeunitwidth}{t}
\measuremaxwidth{\recipedescnaturalwidth}{ground black pepper}
\measuremaxwidth{\recipeamountwidth}{1/4}
\measuremaxwidth{\recipeunitwidth}{t}
\measuremaxwidth{\recipedescnaturalwidth}{white pepper, if desired}
\measuremaxwidth{\recipeamountwidth}{2}
\measuremaxwidth{\recipeunitwidth}{T}
\measuremaxwidth{\recipedescnaturalwidth}{Alstertor Dusseldorf-style mustard}
\measuremaxwidth{\recipeamountwidth}{2}
\measuremaxwidth{\recipeunitwidth}{T}
\measuremaxwidth{\recipedescnaturalwidth}{beef or veg Better-Than-Bouillon}
\measuremaxwidth{\recipeamountwidth}{1/2}
\measuremaxwidth{\recipeunitwidth}{c}
\measuremaxwidth{\recipedescnaturalwidth}{MCT oil}
\measuremaxwidth{\recipeamountwidth}{1}
\measuremaxwidth{\recipeunitwidth}{\mbox{}}
\measuremaxwidth{\recipedescnaturalwidth}{green onion bunch}
\setlength{\recipeingredientdescwidth}{\dimexpr\recipedescnaturalwidth+\recipewidthbuffer\relax}
\setlength{\recipeingredientblockwidth}{\dimexpr\recipeamountwidth+\recipeamountunitgap+\recipeunitwidth+\recipeunitdescgap+\recipeingredientdescwidth\relax}
\ifdim\recipeingredientblockwidth>0.48\textwidth
  \setlength{\recipeingredientblockwidth}{0.48\textwidth}
  \setlength{\recipeingredientdescwidth}{\dimexpr\recipeingredientblockwidth-\recipeamountwidth-\recipeamountunitgap-\recipeunitwidth-\recipeunitdescgap\relax}
\fi
\setlength{\recipestepwidth}{\dimexpr\textwidth-\recipecolumngap-\recipeingredientblockwidth\relax}

{\setlength{\parskip}{0pt}%
\begin{longtable}{@{}>{\raggedright\arraybackslash}p{\recipestepwidth}@{\hspace{\recipecolumngap}}>{\raggedright\arraybackslash}p{\recipeingredientblockwidth}@{}}
\toprule
\endfirsthead
\toprule
\endhead
\bottomrule
\endfoot
\bottomrule
\endlastfoot
Scrub and rinse.& \parbox[t]{\recipeingredientblockwidth}{\begin{tabular}[t]{@{}>{\raggedleft\arraybackslash}p{\recipeamountwidth}@{\hspace{\recipeamountunitgap}}>{\raggedright\arraybackslash}p{\recipeunitwidth}@{\hspace{\recipeunitdescgap}}>{\raggedright\arraybackslash}p{\recipeingredientdescwidth}@{}}%
5&lbs&medium russet potatoes\\
\end{tabular}}\tabularnewline
\midrule
Add 1~cup of water to the Instant Pot. Insert a metal trivet to keep the potatoes above the water. Stack potatoes above the water level. Cook on high pressure for 16~minutes. Allow 10~minutes of natural pressure release.& \parbox[t]{\recipeingredientblockwidth}{\begin{tabular}[t]{@{}>{\raggedleft\arraybackslash}p{\recipeamountwidth}@{\hspace{\recipeamountunitgap}}>{\raggedright\arraybackslash}p{\recipeunitwidth}@{\hspace{\recipeunitdescgap}}>{\raggedright\arraybackslash}p{\recipeingredientdescwidth}@{}}%
1&c&water\\
\end{tabular}}\tabularnewline
\midrule
In a stovetop pot, bring to a boil over medium heat. Once at a rolling boil, remove from heat.& \parbox[t]{\recipeingredientblockwidth}{\begin{tabular}[t]{@{}>{\raggedleft\arraybackslash}p{\recipeamountwidth}@{\hspace{\recipeamountunitgap}}>{\raggedright\arraybackslash}p{\recipeunitwidth}@{\hspace{\recipeunitdescgap}}>{\raggedright\arraybackslash}p{\recipeingredientdescwidth}@{}}%
2&\mbox{}&chopped yellow onions\\
1.5&c&water\\
1/2&c&vinegar\\
2&t&sea salt\\
1&t&ground black pepper\\
1/4&t&white pepper, if desired\\
2&T&Alstertor Dusseldorf-style mustard\\
2&T&beef or veg Better-Than-Bouillon\\
\end{tabular}}\tabularnewline
\midrule
When cool enough to handle, peel potatoes. Cut off and discard any brown spots. Chop into 3/4" to 1" cubes. Collect potato cubes in a 9×12" glass pan, or a (preferably non-metal) bowl. Pour sauce over still-hot potatoes. As they cool, occasionally mix gently. Once cooled, place in the fridge for at least 12~hours. Remove, drizzle MCT oil over potatoes, and toss. As the oil works into the sauced potatoes, it should become creamy and emulsified. If desired, garnish with minced green onions.& \parbox[t]{\recipeingredientblockwidth}{\begin{tabular}[t]{@{}>{\raggedleft\arraybackslash}p{\recipeamountwidth}@{\hspace{\recipeamountunitgap}}>{\raggedright\arraybackslash}p{\recipeunitwidth}@{\hspace{\recipeunitdescgap}}>{\raggedright\arraybackslash}p{\recipeingredientdescwidth}@{}}%
1/2&c&MCT oil\\
1&\mbox{}&green onion bunch\\
\end{tabular}}\tabularnewline
\end{longtable}}

\subsection*{Notes}

\paragrcompact{Variation} Peel and cut 2.5~lbs of russet potatoes into 1/2" cubes up-front.
Cook HP for 6~min and QR pressure.
Use 1~medium onion, 1~cup water, 1/2~cup vinegar, 1~t salt,
1/2~t no-salt, 1/2~t ground black pepper, and 1~T Better-Than-Bouillon.

\end{document}